%; whizzy paragraph -pdf xpdf -latex ./whizzypdfptex.sh
%; whizzy-paragraph "^\\\\begin{frame}\\|\\\\emtext"
% latex beamer presentation.
% platex, latex-beamer でコンパイルすることを想定。 

%     Tokyo Debian Meeting resources
%     Copyright (C) 2012 Junichi Uekawa

%     This program is free software; you can redistribute it and/or modify
%     it under the terms of the GNU General Public License as published by
%     the Free Software Foundation; either version 2 of the License, or
%     (at your option) any later version.

%     This program is distributed in the hope that it will be useful,
%     but WITHOUT ANY WARRANTY; without even the implied warreanty of
%     MERCHANTABILITY or FITNESS FOR A PARTICULAR PURPOSE.  See the
%     GNU General Public License for more details.

%     You should have received a copy of the GNU General Public License
%     along with this program; if not, write to the Free Software
%     Foundation, Inc., 51 Franklin St, Fifth Floor, Boston, MA  02110-1301 USA

\documentclass[cjk,dvipdfmx,12pt]{beamer}
\usetheme{Tokyo}
\usepackage{monthlypresentation}

%  preview (shell-command (concat "evince " (replace-regexp-in-string "tex$" "pdf"(buffer-file-name)) "&")) 
%  presentation (shell-command (concat "xpdf -fullscreen " (replace-regexp-in-string "tex$" "pdf"(buffer-file-name)) "&"))
%  presentation (shell-command (concat "evince " (replace-regexp-in-string "tex$" "pdf"(buffer-file-name)) "&"))

%http://www.naney.org/diki/dk/hyperref.html
%日本語EUC系環境の時
\AtBeginDvi{\special{pdf:tounicode EUC-UCS2}}
%シフトJIS系環境の時
%\AtBeginDvi{\special{pdf:tounicode 90ms-RKSJ-UCS2}}

\newenvironment{commandlinesmall}%
{\VerbatimEnvironment
  \begin{Sbox}\begin{minipage}{1.0\hsize}\begin{fontsize}{8}{8} \begin{BVerbatim}}%
{\end{BVerbatim}\end{fontsize}\end{minipage}\end{Sbox}
  \setlength{\fboxsep}{8pt}
% start on a new paragraph

\vspace{6pt}% skip before
\fcolorbox{dancerdarkblue}{dancerlightblue}{\TheSbox}

\vspace{6pt}% skip after
}
%end of commandlinesmall

\title{東京エリアDebian勉強会}
\subtitle{第107回 2013年12月度}
\author{野島貴英}
\date{2013年12月21日}
\logo{\includegraphics[width=8cm]{image200607/openlogo-light.eps}}

\begin{document}

\begin{frame}
\titlepage{}
\end{frame}

\begin{frame}{設営準備にご協力ください。}
会場設営よろしくおねがいします。
\end{frame}

\begin{frame}{Agenda}
 \begin{minipage}[t]{0.45\hsize}
  \begin{itemize}
   \item 注意事項
	 \begin{itemize}
	  \item 写真はセミナールーム内のみ可です。
          \item 出入りは自由でないので、もし外出したい方は、野島まで一声くださいませ。
	 \end{itemize}
   \item 最近あったDebian関連のイベント報告
	 \begin{itemize}
	  \item 第106回 東京エリアDebian勉強会
	 \end{itemize}
  \end{itemize}
 \end{minipage} 
 \begin{minipage}[t]{0.45\hsize}
  \begin{itemize}
   \item Debian Trivia Quiz
   \item 2013年東京エリアDebian勉強会の振り返り
   \item 2014年計画（BOF形式)
   \item Debian GNU/Hurd 2013
  \end{itemize}
 \end{minipage}
\end{frame}

\section{イベント報告}
\emtext{イベント報告}

\begin{frame}{第106回 東京エリアDebian勉強会}

 東京エリアDebian勉強会106回めは杉並区荻窪で開催されました。
８名の参加者がありました。

\begin{itemize}
\item 野島さんが、debian sidでのwaylandの動かし方、内部構造、周辺技術の紹介を行いました。また、KMS/DRM(i810ドライバ)を使い、X無しにて、実際にwaylandを動かし、プロジェクタに出力してwaylandのデモを行いました。
\item 上川さんが、emacs上で動く、リモート環境のファイル操作に関する拡張のtrampについてデモと説明を行いました。
\end{itemize}

 全体としてハックに活用できるネタが多かった印象でした。
 
 宴会は「はなの舞」にて行いました。


\end{frame}

\section{Debian Trivia Quiz}
\emtext{Debian Trivia Quiz}
\begin{frame}{Debian Trivia Quiz}

  Debian の常識、もちろん知ってますよね?
知らないなんて恥ずかしくて、知らないとは言えないあんなことやこんなこと、
みんなで確認してみましょう。

今回の出題範囲は\url{debian-devel-announce@lists.debian.org},
\url{debian-devel@lists.debian.org} に投稿された
内容などからです。

\end{frame}

\subsection{問題}
 %; whizzy-master ../debianmeetingresume201311.tex
% $B0J>e$N@_Dj$r$7$F$$$k$?$a!"$3$N%U%!%$%k$G(B M-x whizzytex $B$9$k$H!"(Bwhizzytex$B$,MxMQ$G$-$^$9!#(B
%

\santaku
{2013/12/15$B$K=P$?(BDebian wheezy$B$N%"%C%W%G!<%H$N%P!<%8%g%s$O!)(B}
{7.3}
{7.2}
{7.1}
{A}
{debian wheezy$B$r$*;H$$$N3'MM$OAaB.%"%C%W%G!<%H$7$^$7$g$&!*(B}

\santaku
{Debian$B$b;22C$7$F$$$k(BFOSS$B9W8%<T$K=w@-$rA}$d$=$&1?F0$N;v$r$J$s$H$$$&!)(B}
{Encourage Women in Linux}
{Outreach Program For Women}
{IT$B@o;N(B}
{B}
{Outreach Program For Women($BN,$7$F(BOPW)$B$O!"(B\\
\url{http://gnome.org/opw/}
$B$K>pJs$,$"$j$^$9!#(B}

\santaku
{$B@hF|(BDebian$B$N(BTechnical Committee$B$K%a%s%P$,A}$($^$7$?!#$H$3$m$G!"(B
 $B8=:_$N(BTechnical Comittee$B$N(Bchair man$B$OC/$G$7$g$&(B?}
{Takahide Nojima}
{Lucas Nussbaum}
{Bdale Garbee}
{C}
{$B85(BHP$B$N(BLinux$BItLg$N(BCTO$B$rL3$a$?$H$$$&7PNr$N;}$A<g$G$9!#(Bdebconf$B$H$+$K(B
$B;22C(B/$B%S%G%*;kD0$H$+$9$k$H$o$+$k$N$G$9$,!">oO"$5$s$N$h$&$G$9!#(B}



\section{2013年東京エリアDebian勉強会の振り返り}
\emtext{2013年東京エリアDebian勉強会の振り返り}


\section{2014年計画}
\emtext{2014年計画}

{\footnotesize
 \begin{prework}{ 河本拓 }

 インタビュー形式で、参加者のみなさんが「オープンソースに関わったきっかけ」などのお話が聞けたら嬉しいです。2014年最初の勉強会だと思うので、新年の抱負などと合わせて...(?)
\\
追記:初参加です。Linux初心者ですがよろしくお願いします。
\end{prework}

\begin{prework}{ dictoss(杉本　典充) }

基本的にはハックカフェのようにdebianに関する作業をする場として集まり、質問や議論があれば集まった人たちの中で話す、という場を提供することにするとよさそう。
参加の条件として行った作業や成した事、成せなかった事をまとめてブログ記事を必ず１つその場で書くこととする、というのはどうだろうか。そのとき各個人のブログサイトではなく、「勉強会ブログ」のようなサイトで一元的に記事を登録するようにして、対外的にdeibanのトレンドを発信することも兼ねるというのはどうだろうか。
ただ参加者の間でパッケージやdebianの仕組みについて勉強会が必要と判断することもあるので、その場合はセミナー形式で開催すればいいと思う。
\end{prework}

\begin{prework}{ 吉野(yy\_{}y\_{}ja\_{}jp) }

特に変わらないと思っていました
\end{prework}

\begin{prework}{ 野島　貴英 }

原稿／プレゼン集まればセミナー形式、集まらなければ、ハッカソン（作業時間）にするのがよいかなー？とりあえず、来年こそは、debian組み込みネタか、ハードウェアネタをやりたい。
\end{prework}

}

\section{debian GNU/Hurd 2013}
\emtext{debian GNU/Hurd 2013}

\section{今後のイベント}
\emtext{今後のイベント}
\begin{frame}{今後のイベント}
\begin{itemize}
 \item 2014年1月 Debian勉強会
\end{itemize}
\end{frame}

\section{今日の宴会場所}
\emtext{今日の宴会場所}
\begin{frame}{今日の宴会場所}
未定
\end{frame}

\end{document}

;;; Local Variables: ***
;;; outline-regexp: "\\([ 	]*\\\\\\(documentstyle\\|documentclass\\|emtext\\|section\\|begin{frame}\\)\\*?[ 	]*[[{]\\|[]+\\)" ***
;;; End: ***
